\chapter{Results}\label{ch:results}

This chapter reports the results from the evaluation methods in Chapter~\ref{ch:evaluation}. The results are grouped by the same evidence chain used in the thesis: sensor agreement with Shimmer, pilot field-worn trend agreement, detection performance on the labelled train-sitting dataset, and embedded timing of the compact KNN.

\section{Sensor Agreement Results}\label{sec:results_sensor_agreement}

\subsection{Co-Located Shimmer/XIAO Agreement}

Figure~\ref{fig:shimmer_xiao_comparison} shows the aligned time-series comparison for the co-located recording. The XIAO and Shimmer signals show the same main movement peaks across all acceleration and gyroscope channels. Across all axes, acceleration correlations ranged from 0.856 to 0.923 and gyroscope correlations from 0.816 to 0.864 (Table~\ref{tab:colocated_validation_results}). The estimated lag stayed between \(-0.020\) and 0~s, and NRMSE values ranged from 0.425 to 0.613.

\begin{figure*}[!tbp]
    \centering
    \includegraphics[width=0.94\textwidth,height=0.58\textheight,keepaspectratio]{shimmer_xiao_comparison.png}
    \caption{Aligned co-located waveform comparison between the Shimmer reference and the custom XIAO-based wearable, across the three acceleration and three gyroscope axes.}
    \label{fig:shimmer_xiao_comparison}
\end{figure*}

\begin{table}[H]
    \centering
    \scriptsize
    \caption{Co-located agreement after alignment.}
    \label{tab:colocated_validation_results}
    \begin{tabular*}{\linewidth}{@{\extracolsep{\fill}}llrrrr@{}}
        \toprule
        Signal & Axis & r & Max r & Lag (s) & NRMSE \\
        \midrule
        Accel & X & 0.907 & 0.910 & -0.010 & 0.517 \\
        Accel & Y & 0.856 & 0.856 & -0.010 & 0.613 \\
        Accel & Z & 0.923 & 0.923 & -0.010 & 0.425 \\
        Gyro  & X & 0.838 & 0.838 &  0.000 & 0.568 \\
        Gyro  & Y & 0.816 & 0.816 &  0.000 & 0.601 \\
        Gyro  & Z & 0.864 & 0.881 & -0.020 & 0.521 \\
        \bottomrule
    \end{tabular*}
\end{table}

The co-located results support the use of the XIAO signal for trend-level movement detection. The correlations indicate similar temporal patterns, while the non-zero NRMSE values show that the two sensors should not be treated as amplitude-equivalent measurement systems.

\subsection{Pilot Field-Worn Trend Agreement}

The field-worn trend comparison used the two usable paired recordings from the pilot study. Acceleration motion magnitude correlated at 0.918 for P1 and 0.855 for P2, while gyroscope magnitude correlated at 0.908 for P1 and 0.864 for P2. After lag correction, maximum cross-correlations reached 0.915--0.981. The estimated best lag was 1.2--1.3~s (Table~\ref{tab:field_trend_results}). Figure~\ref{fig:field_trend_comparison} shows the corresponding long trend comparison.

\begin{table}[H]
    \centering
    \scriptsize
    \caption{Field-worn trend agreement in the pilot study. P1 and P2 are participant identifiers.}
    \label{tab:field_trend_results}
    \begin{tabular*}{\linewidth}{@{\extracolsep{\fill}}llrrrr@{}}
        \toprule
        Part. & Signal & r & Max r & Lag (s) & NRMSE \\
        \midrule
        P1 & Accel.\ mag. & 0.918 & 0.938 & 1.20 & 0.530 \\
        P1 & Gyro mag.    & 0.908 & 0.974 & 1.20 & 0.232 \\
        P2 & Accel.\ mag. & 0.855 & 0.915 & 1.30 & 0.476 \\
        P2 & Gyro mag.    & 0.864 & 0.981 & 1.30 & 0.206 \\
        \bottomrule
    \end{tabular*}
\end{table}

\begin{figure*}[!tbp]
    \centering
    \includegraphics[width=0.94\textwidth,height=0.42\textheight,keepaspectratio]{c5_field_trend_comparison.png}
    \caption{Long field-worn trend comparison for two pilot-study participants. Signals are shown as acceleration motion magnitude and gyroscope magnitude over the overlapping recording period, after 10~Hz resampling and centred 5~s moving-average smoothing for visualisation only.}
    \label{fig:field_trend_comparison}
\end{figure*}

The field-worn correlations are lower than a controlled laboratory equivalence test would require, but that is not the claim being made here. The result shows that the XIAO captured long movement trends similar to the Shimmer reference during the pilot outing. The 1.2--1.3~s lag indicates a timing-chain offset in the field setup and is discussed in Chapter~\ref{ch:discussion}.

\section{Detection Results}\label{sec:results_detection}

\subsection{Classifier Comparison}

Table~\ref{tab:ch6_model_results} reports the held-out performance of all five classifiers on participant 5. Static sitting outnumbers pelvic rotation 21:1 in the test set, so balanced accuracy and pelvic-rotation F1-score are more informative than accuracy alone. KNN achieved the highest F1-score at 0.957.

\begin{table}[H]
    \centering
    \footnotesize
    \caption{Held-out participant-5 classification results. Precision, recall, and F1 are for the pelvic-rotation class.}
    \label{tab:ch6_model_results}
    \begin{tabular*}{\linewidth}{@{\extracolsep{\fill}}lrrrrr@{}}
        \toprule
        Model & Acc & Bal.\ acc & P & R & F1 \\
        \midrule
        LR  & 0.909 & 0.912 & 0.320 & 0.915 & 0.474 \\
        SVM & 0.964 & 0.970 & 0.557 & 0.976 & 0.709 \\
        DT  & 0.978 & 0.925 & 0.711 & 0.867 & 0.781 \\
        RF  & 0.993 & 0.921 & 0.993 & 0.842 & 0.911 \\
        KNN & 0.996 & 0.969 & 0.975 & 0.939 & 0.957 \\
        \bottomrule
    \end{tabular*}
\end{table}

Out of 165 pelvic-rotation windows, the selected KNN correctly identified 155 and missed 10. Out of 3{,}497 static-sitting windows, it classified 3{,}493 correctly and misclassified 4 as pelvic rotation. The corresponding confusion matrix is shown in Figure~\ref{fig:ch6_knn_confusion_matrix}.

\begin{figure}[!tbp]
    \centering
    \includegraphics[width=0.76\linewidth]{ch6_person5_confusion_matrix_best_f1.png}
    \caption{Confusion matrix of the selected \(k=7\) KNN on held-out participant 5.}
    \label{fig:ch6_knn_confusion_matrix}
\end{figure}

\subsection{Compact KNN and On-Device Timing}

The compact prototype-based KNN (12 features, 128 KMeans prototypes), described in Section~\ref{sec:compact_knn}, achieved accuracy 0.996, balanced accuracy 0.963, precision 0.994, recall 0.927, and an F1-score of 0.959 on the held-out participant. The full baseline used 90 features and 9{,}391 reference windows, whereas the compact model used 12 features and 128 prototypes (Figure~\ref{fig:full_vs_compact_knn}). The compact model therefore had slightly higher precision and lower recall than the full KNN, and the comparison changes both the feature set and the reference-set size.

\begin{figure*}[!tbp]
    \centering
    \includegraphics[width=0.92\textwidth]{ch6_full_vs_compact_knn_comparison.png}
    \caption{Comparison between the offline full KNN (90 features, 9{,}391 reference windows) and the compact prototype-based KNN (12 features, 128 KMeans prototypes) used for embedded deployment.}
    \label{fig:full_vs_compact_knn}
\end{figure*}

Table~\ref{tab:knn_timing} summarises per-window inference timing measured across 24 consecutive non-overlapping 2~s windows on the XIAO nRF54L15 Sense. The timing measurement includes the compact feature extraction path and the KNN distance computation for the stored prototypes.

\begin{table}[H]
    \centering
    \footnotesize
    \caption{On-device execution time of the compact KNN.}
    \label{tab:knn_timing}
    \begin{tabularx}{\linewidth}{lXX}
        \toprule
        Stage & Mean & Range \\
        \midrule
        KNN inference only           & 2.70~ms & 2.67--2.73~ms \\
        Features + KNN total         & 2.73~ms & 2.70--2.76~ms \\
        \bottomrule
    \end{tabularx}
\end{table}

The 2.73~ms total is below 0.15\% of the 2{,}000~ms window period. This percentage describes compute time only; it is not a power-consumption or battery-life measurement. The deployed firmware uses 56.3~KB of flash and 9.3~KB of RAM, leaving headroom within the target's available flash and RAM budget for the Zephyr RTOS, BLE stack, and future extensions.
